% ===========================================================================
% MCM/ICM Mathematical Contest in Modeling LaTeX Template
% English paper format
%
% Format requirements:
%   - Summary on separate page, max 1 page
%   - Three-line tables (booktabs: toprule/midrule/bottomrule, 1.5pt, no vertical lines)
%   - Table captions above, figure captions below
%   - Figures at 600dpi, black-and-white distinguishable (line styles + legends)
%   - Formulas centered, right-aligned numbering (only when referenced)
%   - References 5-10, standard academic format
%   - Times New Roman body + math
%   - Margins 2.5cm
%   - Body no more than 20 pages
% ===========================================================================

\documentclass[12pt,a4paper]{article}

% ---------- Page setup: 2.5cm margins ----------
\usepackage[top=2.5cm,bottom=2.5cm,left=2.5cm,right=2.5cm]{geometry}

% ---------- Font: Times New Roman ----------
\usepackage{times}
\usepackage[T1]{fontenc}

% ---------- Three-line tables (booktabs, 1.5pt) ----------
\usepackage{booktabs}
\renewcommand{\toprule}{\specialrule{1.5pt}{0pt}{0pt}}
\renewcommand{\midrule}{\specialrule{1.5pt}{0pt}{0pt}}
\renewcommand{\bottomrule}{\specialrule{1.5pt}{0pt}{0pt}}

% ---------- Mathematics ----------
\usepackage{amsmath,amssymb}

% ---------- Graphics ----------
\usepackage{graphicx}
\usepackage[font=small,labelfont=bf,labelsep=space]{caption}
% Figure captions below figures (default)
\captionsetup[figure]{position=below}
% Table captions above tables
\captionsetup[table]{position=above}

% ---------- Hyperlinks ----------
\usepackage[colorlinks=true,linkcolor=black,citecolor=black,urlcolor=black]{hyperref}

% ---------- References ----------
\usepackage[numbers,sort&compress]{natbib}

% ---------- Line spacing: 1.5 ----------
\usepackage{setspace}
\onehalfspacing

% ---------- Appendix ----------
\usepackage[titletoc]{appendix}

% ---------- Code listings ----------
\usepackage{listings}
\usepackage{xcolor}
\lstset{
    basicstyle=\ttfamily\small,
    breaklines=true,
    frame=single,
    numbers=left,
    numberstyle=\tiny,
}

% ===========================================================================
% Paper information (fill in here)
% ===========================================================================
\newcommand{\mcmtitle}{Paper Title}
\newcommand{\mcmkeywords}{keyword1; keyword2; keyword3; keyword4}

% ===========================================================================
% Document begins
% ===========================================================================
\begin{document}

% ===================================================================
% Summary page (separate page)
% ===================================================================
\begin{center}
    {\Large\bfseries \mcmtitle}

    \vspace{18pt}
    {\large\bfseries Summary}
\end{center}

\vspace{6pt}

% --- Notes ---
% Summary requirements:
%   1. Five elements: restatement / approach / model / results / conclusions
%   2. 3-5 keywords
%   3. No formulas, tables, or figures
%   4. No more than 1 page

This paper addresses the problem of \underline{\hspace{4cm}} by developing
a \underline{\hspace{2cm}} model. First, \underline{\hspace{4cm}}.
Then, \underline{\hspace{4cm}}. Finally, \underline{\hspace{4cm}}.

The main results are as follows:
\begin{enumerate}
    \item Problem 1: \underline{\hspace{6cm}}
    \item Problem 2: \underline{\hspace{6cm}}
    \item Problem 3: \underline{\hspace{6cm}}
\end{enumerate}

\vspace{12pt}
\noindent {\bfseries Keywords:} \mcmkeywords

\newpage

% ===================================================================
% Body begins here (max 20 pages)
% ===================================================================

% ===================================================================
% 1. Introduction
% ===================================================================
\section{Introduction}

\subsection{Background}
% Problem background

\subsection{Problem Statement}
% Specific problem statement

\subsection{Objectives}
% List of objectives

% ===================================================================
% 2. Assumptions and Justifications
% ===================================================================
\section{Assumptions and Justifications}

\subsection{Basic Assumptions}
\begin{itemize}
    \item Data provided by the problem is authentic and reliable.
    \item \underline{\hspace{6cm}}
    \item \underline{\hspace{6cm}}
\end{itemize}

\subsection{Notation}
% Use three-line table for notation
\begin{table}[htbp]
    \centering
    \caption{Notation and Symbols}
    \label{tab:symbols}
    \begin{tabular}{ccc}
        \toprule
        Symbol & Description & Unit \\
        \midrule
        $x$ & \underline{\hspace{3cm}} & \underline{\hspace{2cm}} \\
        $y$ & \underline{\hspace{3cm}} & \underline{\hspace{2cm}} \\
        \bottomrule
    \end{tabular}
\end{table}

% ===================================================================
% 3. Model Development
% ===================================================================
\section{Model Development}

\subsection{Model 1: \underline{\hspace{3cm}}}
% Model construction
% Formulas (centered, numbered only when referenced)
% Solution method
% Results
% Results in three-line tables (caption above)

\begin{table}[htbp]
    \centering
    \caption{Model 1 Results}
    \label{tab:model1}
    \begin{tabular}{ccc}
        \toprule
        Metric & Value & Notes \\
        \midrule
        \underline{\hspace{2cm}} & \underline{\hspace{2cm}} & \underline{\hspace{2cm}} \\
        \bottomrule
    \end{tabular}
\end{table}

% Example formula (numbered only when referenced)
\begin{equation}
    y = f(x) = ax + b
    \label{eq:linear}
\end{equation}

\subsection{Model 2: \underline{\hspace{3cm}}}
% Model construction, solution, results

% ===================================================================
% 4. Model Evaluation
% ===================================================================
\section{Model Evaluation}

\subsection{Strengths}
\subsection{Weaknesses}
\subsection{Potential Improvements}

% ===================================================================
% 5. Sensitivity Analysis
% ===================================================================
\section{Sensitivity Analysis}

% ===================================================================
% 6. Conclusions
% ===================================================================
\section{Conclusions}

% ===================================================================
% References (5-10 entries)
% Standard academic format
% ===================================================================
\begin{thebibliography}{99}

\bibitem{ref1} Author A, Author B. Article Title[J]. \textit{Journal Name}, Year, Volume(Issue): Pages.

\bibitem{ref2} Author C. \textit{Book Title}[M]. Publisher Location: Publisher, Year.

\bibitem{ref3} Author D. \textit{Dissertation Title}[D]. University, Year.

\bibitem{ref4} Author E. Document Title[EB/OL]. URL, Access Date.

\bibitem{ref5} Author F, Author G. Article Title[J]. \textit{Journal Name}, Year, Volume(Issue): Pages.

\end{thebibliography}

% ===================================================================
% Appendix (separate page, include all runnable code)
% ===================================================================
\newpage
\appendix
\section{Appendix A: Program Code}

% ------ Example code ------
% \begin{lstlisting}[caption={Main Program},label=lst:main]
% import numpy as np
% import pandas as pd
% # ... your code here
% \end{lstlisting}

\end{document}
